\documentclass{standalone}
\usepackage{tikz}
\usepackage{pgfplots}
\usetikzlibrary{patterns,decorations.markings,backgrounds,arrows.meta}

\usepackage{newtxtext,newtxmath}
\usepgfplotslibrary{colorbrewer,colormaps}

\pgfplotsset{compat=1.18}
\pgfplotsset{colormap/RdBu}

\def\image{../pngs/d2.25_w35white.png}
\def\secondimage{../pngs/d2.25_w35whiteunscaled.png}

\def\figurewidth{4.46}
\def\pagewidth{13.5}
\def\geomscale{0.3670781893}

\def\xfirstedge{-5.6}
\def\yfirstedge{-1.42}
\def\xsecondedge{-1.57}
\def\ysecondedge{-2.5}
\def\xmin{-6.075}
\def\xmax{6.075}
\def\ymax{4.6}
\def\ymin{-8}

\def\ysmallgap{3.37}
\def\xsmallgap{-5.97}
\def\depth{0.33}
\def\width{5.28}

\newcommand{\fsize}{\fontsize{8.7}{10.44}\selectfont}

\pgfmathsetlengthmacro{\colorbarwidth}
{0.33*\pagewidth cm}

\pgfmathsetlengthmacro{\colorbarheight}
{0.08*0.33*\pagewidth cm}

\pgfmathsetlengthmacro{\colorbarlinewidth}
{0.0020*\pagewidth cm}

\pgfmathsetlengthmacro{\colorbarticklength}
{0.0075*\pagewidth cm}

\begin{document}

\begin{tikzpicture}[x=\geomscale cm,y=\geomscale cm]

	\path[use as bounding box]
	(-1,-0.7) rectangle (12.9,10.6);

	\begin{scope}[shift={(6.075,6)}]

		\clip (\xmin,\ymin) rectangle (\xmax,\ymax);

		\node[inner sep=0pt] (gap) at (0,0)
		{\includegraphics[width=\figurewidth cm]{\image}};

		\draw[thick]
		(\xmin,\yfirstedge)
		-- (\xfirstedge,\yfirstedge)
		-- (\xfirstedge,\ysecondedge)
		-- (\xsecondedge,\ysecondedge)
		-- (\xsecondedge,\yfirstedge)
		-- (\xmax,\yfirstedge);

		\node[
			inner sep=0pt,
			yshift=0.47cm
		] (gap2) at (gap.north)
		{\includegraphics[
				trim=0cm 31cm 0cm 20cm,
				clip,
				width=\figurewidth cm
			]{\secondimage}};

		\draw[very thin]
		(\xmin,\ysmallgap)
		-- (\xsmallgap,\ysmallgap)
		-- (\xsmallgap,\ysmallgap-\depth)
		-- (\xsmallgap+\width,\ysmallgap-\depth)
		-- (\xsmallgap+\width,\ysmallgap)
		-- (\xmax,\ysmallgap);

		\def\lengthWall{0.5}

		\foreach \x in {-5.7,-1.1,-0.6,...,6.0} {
			\draw
			(\x,\yfirstedge)
			-- ++(225:\lengthWall);
		}

		\foreach \x in {-5.6,-5.2,...,-1.47} {
			\draw
			(\x,\ysecondedge)
			-- ++(225:\lengthWall);
		}

		\draw
		(\xfirstedge,-1.7)
		-- ++(225:\lengthWall);

		\draw
		(\xfirstedge,-2.1)
		-- ++(225:\lengthWall);

		\def\secondlengthWall{0.1}

		\foreach \x in {-0.6,-0.5,...,6.1} {
			\draw[ultra thin]
			(\x,\ysmallgap)
			-- ++(225:\secondlengthWall);
		}

		\foreach \x in {-5.97,-5.87,...,-0.7} {
			\draw[ultra thin]
			(\x,\ysmallgap-\depth)
			-- ++(225:\secondlengthWall);
		}

		\draw[ultra thin]
		(\xsmallgap,\ysmallgap-\depth/7)
		-- ++(225:\secondlengthWall);

		\draw[ultra thin]
		(\xsmallgap,\ysmallgap-2*\depth/5)
		-- ++(225:\secondlengthWall);

		\draw[ultra thin]
		(\xsmallgap,\ysmallgap-3.5*\depth/5)
		-- ++(225:\secondlengthWall);

		\draw[ultra thin]
		(\xmin+0.05,\ysmallgap)
		-- ++(225:\secondlengthWall);

	\end{scope}

	% Colorbar geometry is referenced to a 13.5 cm page width.
	% The actual image itself is 4.46 cm wide.

	\coordinate (colorbarcenter) at (6.075,3);

	\pgfplotscolorbardrawstandalone[
		colormap={example}{
			indices of colormap=(
				10,9,8,7,6,5,4,3,2,1,0 of RdBu
			),
		},
		point meta min=-0.08,
		point meta max=0.08,
		colorbar horizontal,
		samples=5,
		scaled x ticks=false,
		colorbar style={
			xtick={-0.08,-0.04,...,0.08},
			xticklabel style={
				/pgf/number format/fixed,
				/pgf/number format/precision=3,
				font=\fsize,
			},
			xlabel={$v$},
			xlabel style={
				font=\fsize,
				at={(axis description cs:0.5,-1.2)},
			},
			at={(colorbarcenter)},
			anchor=north,
			width=\colorbarwidth,
			height=\colorbarheight,
			xtick pos=lower,
			xtick style={
				black,
				line width=\colorbarlinewidth,
			},
			tickwidth=\colorbarticklength,
			axis line style={
				line width=\colorbarlinewidth,
			},
		},
	]

\end{tikzpicture}

\end{document}
